\documentclass{article}
\usepackage[T1]{fontenc}
\usepackage{lmodern}
\usepackage{graphicx}
\usepackage{amsmath,amssymb}
\usepackage[a4paper,margin=2cm]{geometry}
\usepackage[absolute,overlay]{textpos}
\setlength{\TPHorizModule}{1mm}
\setlength{\TPVertModule}{1mm}
\usepackage[indonesian]{babel}
\usepackage{hyperref}
\setlength{\emergencystretch}{2em}
% unit: Halaman 13

\begin{document}

% === Halaman 13 (python/native) ===

Contoh pengiriman pesan 1: 

• Misalkan Bob mengirim pesan m = 16 kepada Alice. Bob mengenkripsi m dengan 

kunci publik Alice (e, n) = (79, 3337) sebagai berikut:

           c = me mod n  = 1679 mod 3337 = 1493

   Bob mengirim c = 1493 kepada Alice

• Alice mendekripsi cipherteks c = 1493 dengan kunci privat miliknya, yaitu (d, n) = 

(1019, 3337) sebagai berikut:


 m = cd mod n  = 14931019 mod 3337 = 16

   Alice mendapatkan kembali plainteks dari Bob, yaitu m = 16

13

\end{document}
